%-*-latex-*-
%-*-latex-*-
\newcommand\COURSE{ciss350}
\newcommand\ASSESSMENT{q10203}
\newcommand\ASSESSMENTTYPE{Quiz}
\newcommand\POINTS{\textwhite{xxx/xxx}}

\input{myquizpreamble}
\input{yliow}
\input{\COURSE}
%-*-latex-*-

%-*-latex-*-

%-*-latex-*-
\renewcommand\TITLE{\ASSESSMENTTYPE \ \ASSESSMENT}

\renewcommand\EMAIL{}
\renewcommand\AUTHOR{eagraham1@cougars.ccis.edu}

\textwidth=6in
\begin{document}
\topmatterthree

\LATEXHELPTHREEFIVEZERO

\TURNPAGE

%------------------------------------------------------------------------------
\newpage
\textsc{Runtime for linear filter on 1D array }

The goal is to compute the big-O runtime of an array averaging computation
for an array of squares (as well as the initialization of the array of
squares):
\begin{Verbatim}[frame=single,fontsize=\small]
for i = 0, ..., n - 1:
    y[i] = i * i

for i = 1, ..., n - 2:
    z[i] = (y[i - 1] + y[i] + y[i + 1]) / 3

z[0] = z[1]
z[n - 1] = z[n - 2]
\end{Verbatim}
The array \verb!y! contains a collection of (some) values, which in this
case is an array of squares.
Recall from CISS240 that this is an example of a running average computation,
each \verb!z[i]! is the average of 3 values of \verb!y! around the
same index \verb!i!.
(We say the averaging window has size 3.)
This is a basic statistical computation that can be used for instance in
smoothing out a graph of stock prices throughout the day.
In signal processing, this is also an example of one of the simplest
linear filters.
Another example of a linear filter is one where
\verb!z[i] = (y[i - 1] + 2 * y[i] + y[i + 1]) / 4!
where the central value has a higher weight;
in this case the weights are $(1, 2, 1)$.
And in the case of image processing the arrays would be
2D arrays of pixels and the average calculation would be averaging over
9 pixels around a pixel location \verb!(r, c)! (including itself)
and the computation would remove \lq\lq noise" from the image.
An example is the Gaussian filter
\begin{Verbatim}[fontsize=\small]
z[r][c] = (    y[r - 1][c - 1] + 2 * y[r - 1][c] +     y[r - 1][c + 1]
          (2 * y[  r  ][c - 1] + 4 * y[  r  ][c] + 2 * y[  r  ][c + 1]
          (    y[r + 1][c - 1] + 2 * y[r + 1][c] +     y[r + 1][c + 1]) / 16
\end{Verbatim}

Back to our algorithm,
take note of this: we are running two for-loops.
In real wallclock time, the total runtime would take more time than the time
taken for each individual for-loop.
But what about the \underline{\textit{big-O}} of the total runtime?

%------------------------------------------------------------------------------
\nextq
Rewrite the above algorithm with
goto statements, assign constant times to each statement and
compute the time taken as a function of $n$, $t_1$, $t_2$, $t_3$, ....
where $t_1, t_2, t_3, ...$ are runtimes for each statement in your
algorithm. Make sure you
number these times starting with $t_1$, then
$t_2$, etc.
Include the number of times each statement is executed.
\\
\ANSWER
\begin{answercode}
          i = 0                                      t1      1
LOOP1:    if i >= n:                                 t2      n + 1
              goto ENDLOOP1                          t3      1
          y[i] = i * i                               t4      n
          i = i + 1                                  t5      n
          goto LOOP1                                 t6      n
          ENDLOOP1

          i = 1                                      t7      1
LOOP2:    if i >= n:                                 t8      n - 1
               goto ENDLOOP2                         t9      1
          z[i] = y(y[i - 1] + y[i] + y[i + 1] / 3    t10     n - 2
          i = i + 1                                  t11     n - 2
          goto LOOP2                                 t12     n - 2
          ENDLOOP2

          z[0] = z[1]                                t13     1
          z[n - 1] = z[n - 2]                        t14     1
          
\end{answercode}

%------------------------------------------------------------------------------
\nextq
Let $T(n)$ be the runtime function of the above algorithm.
Write $T(n)$ as a polynomial in descending power of $n$, i.e.,
the \textit{highest degree term appearing first}.
\\
\ANSWER
\begin{answerlong}
$T(n) = (t_{2} + t_{4} + t_{5} + t_{6} + t_{8} + t_{10} + t_{11} + t_{12})n + (t_{1} + t_{2} + t_{3} + t_{7} - t_{8} + t_{9} - 2t_{10} - 2t_{11} - 2t_{12} + t_{13} + t_{14})$
\end{answerlong}
(Correct the above.)

%------------------------------------------------------------------------------
\nextq
Simplify the runtime function $T(n)$
by giving names $A$, $B$, $C$, $D$, ...
to the constants of the function from (a). (The $t_1, t_2, ...$ should
disappear.)
\\
\ANSWER
\begin{answerlong}
$T(n) = A n + B$
\end{answerlong}
(Correct the above.)

%------------------------------------------------------------------------------
\nextq
Fudge away the constants and write down the simplest $g(n)$ such
that the runtime is a big-O of your $g(n)$.
Your $g(n)$ should be either $n$ or $n^2$ or $n^3$ or $n^{102.25}$ ...
\\
\ANSWER
\begin{answerlong}
$T(n) = O(n)$
\end{answerlong}
(Correct the above.)
\qed

%------------------------------------------------------------------------------
%-*-latex-*-
\newpage
\input{instructions}
\end{document}

